\chapter{Synthesis, limitations and future work}
\label{ch:synthesis}

\section{A target-first synthesis}
\label{sec:synthesis-target-first}

The four research chapters approach distributional inference from different
directions. Chapter~\ref{ch:research-a} makes dynamic quantile methodology
usable through a software and computational workflow. Chapter
\ref{ch:research-b} uses that framework to correct and synthesize heterogeneous
hydrologic products. Chapter~\ref{ch:research-c} replaces linear state
evolution for the quantile location with fixed nonlinear recurrent features.
Chapter~\ref{ch:research-d} changes the inferential object from pointwise
conditional quantiles to fixed-content interval placement and tolerance
actions. These are linked contributions, but they are not four implementations
of one model.

Table~\ref{tab:synthesis-comparison} compares the projects along the dimensions
that determine their statistical interpretations. The target and update
columns are especially important. Chapters~\ref{ch:research-a}--
\ref{ch:research-c} use AL or exAL likelihood contributions to study
conditional quantiles and, in selected analyses, synthesized future-response
distributions. Chapter~\ref{ch:research-d} uses a residual-product loss to
define interval endpoints and a separate repeated-sampling calibration to
construct a tolerance action. A common latent-variable algebra does not erase
this difference.

\begin{table}[htbp]
\centering
\begingroup
\small
\setlength{\tabcolsep}{3pt}
\renewcommand{\arraystretch}{1.08}
\begin{tabularx}{\textwidth}{@{}p{0.14\textwidth}YYYY@{}}
\toprule
\textbf{Chapter} & \textbf{Primary target} & \textbf{Dependence representation}
& \textbf{Update and computation} & \textbf{Principal evidence} \\
\midrule
2: \pkg{exdqlm} & Dynamic or static conditional quantile; post-fit predictive
synthesis & Linear latent states, transfer components, or static regression &
AL/exAL likelihood; MCMC, LDVB, and ISVB interfaces & Software examples,
diagnostics, held-out scores, and exact version contract \\
3: hydrology & Source-corrected conditional quantiles and a synthesized
future-response distribution & Shared latent quantile process with source-
specific discrepancies and transfer covariates & exAL dynamic model with LDVB
and supporting MCMC algorithms & Five forecast origins; separate 28-day and
common 8-day comparisons \\
4: Q--DESN & Single or joint conditional quantiles & Fixed deep-reservoir
features followed by Bayesian regression & AL/exAL working likelihood; MCMC
and model-specific VB--LD & Known-DGP simulations and retrospective GloFAS and
PriceFM applications \\
5: MTI & Fixed-content interval placement; calibrated empirical tolerance
action & Static endpoints or proposed regression and dynamic endpoint states &
Generalized Bayes for fixed targets; numerical scan calibration for TCSP &
Population results, simulation validation, and a descriptive pharmaceutical
application \\
\bottomrule
\end{tabularx}
\endgroup
\caption{Comparison of the four research chapters. The entries identify the
main inferential object and the strongest evidence available; they do not imply
that the validation designs are interchangeable.}
\label{tab:synthesis-comparison}
\end{table}

Three connections emerge from this comparison. First, distributional targets
can be modeled without committing every analysis to a complete parametric
response distribution. The quantile chapters target selected conditional
features, while the MTI chapter begins from the geometry of a fixed-content
interval. Second, temporal structure can enter through a latent linear state,
a fixed nonlinear recurrent feature map, or evolving endpoint functions. The
interpretation of the state depends on the target it indexes. Third,
computational tractability is achieved by exposing conditionally Gaussian
structure, but posterior simulation, variational approximation, and
generalized-posterior computation retain different inferential meanings.

\section{What the four projects establish}
\label{sec:synthesis-findings}

\subsection{From methodology to a usable dynamic quantile workflow}

Chapter~\ref{ch:research-a} shows how exAL and exDQLM foundations can be
organized as a coherent analysis workflow. The \pkg{exdqlm} interface connects
model construction, posterior computation, state and component summaries,
forecasting, calibration diagnostics, predictive scoring, and post-fit
quantile synthesis. Its examples establish functionality across dynamic and
static analyses rather than a universal predictive advantage. The Big Tree
study, for example, gives mixed fitted and held-out comparisons among the
baseline, direct-regression, and transfer specifications. The static sparse
example likewise shows the intended computational tradeoff: LDVB is faster in
the reported fits, whereas MCMC gives the smaller held-out quantile error in
all three quantile-specific comparisons and remains the reference when tail or
interval behavior is central.

The chapter also makes provenance part of the statistical contract. The
hydrologic application uses the recorded version 1.1.0 interface, the software
article is fixed to 1.1.1, and the later 1.1.2 repository state is retained as
a support snapshot rather than substituted retroactively. This version
separation is necessary because a package name alone does not identify the
algorithm, default, or fitted object used in an analysis.

\subsection{Source-aware hydrologic synthesis}

Chapter~\ref{ch:research-b} provides the most direct link between the software
and an applied forecasting problem. Its principal empirical finding is that the
selected multivariate exAL specification with an active forecast-window
transfer component attains the lowest 28-day CRPS among the compared Bayesian
variants at each of the five retained origins. The shorter common eight-day
comparison against NWS products is deliberately reported separately and gives
more heterogeneous results. Thus the evidence supports source-aware correction
and synthesis in the studied 28-day design without implying improvement over
every comparator at every horizon.

This application also clarifies what heterogeneous-source modeling requires.
Observations, retrospective products, issued ensembles, and exogenous forecast
summaries have different temporal supports and roles. Treating them as named
source channels allows discrepancies learned before an origin to be propagated
into a forecast-window analysis. It does not remove uncertainty about the
inputs themselves: the future precipitation and soil-water ensembles enter
through deterministic summaries in the current implementation. A fuller
operational analysis would have to propagate that uncertainty and evaluate many
more version-consistent origins.

\subsection{Nonlinear recurrent features for conditional quantiles}

Chapter~\ref{ch:research-c} shows that a fixed deep echo-state reservoir can
serve as a nonlinear design for Bayesian conditional-quantile regression. The
single-level models allow AL or exAL working likelihoods and shrinkage priors;
the joint formulation shares information across an ordered quantile grid
through shrinkage on adjacent coefficient differences. The joint prior can
reduce crossing but does not impose monotonicity, so crossing diagnostics and
post-fit monotone correction retain an explicit role.

The simulations document recovery and forecasting under stated data-generating
mechanisms and selected reservoirs. They do not integrate uncertainty over
reservoir seeds or provide a universal MCMC--VB approximation comparison. The
retrospective applications add two distinct forms of evidence. The GloFAS
analysis produced an ordered issued-window grid with the reported empirical
coverage at one origin, but it used retrospectively observed covariates and the
fit reached its iteration cap before the strict outer tolerance. The PriceFM
comparison froze one selected Q--DESN specification per region across three
held-out folds; its aggregate score is slightly worse than the matched PriceFM
score, with differences that vary by region and fold. These results demonstrate
how the method behaves under the recorded protocols while resisting a broad
ranking from two bounded applications.

\subsection{Fixed-content placement and tolerance inference}

Chapter~\ref{ch:research-d} establishes that fixed probability content leaves
an interval-placement problem. Under its assumptions, the residual-product
criterion identifies a mean-preserving interval, and the mean-tilt parameter
indexes admissible contiguous windows of the prescribed content. Equal-tailed
and shortest intervals appear as distribution-dependent locations on this
path rather than as universal choices of one fixed tilt.

The tolerance contribution is intentionally separate from this population
geometry. TCSP numerically calibrates a retained count from a uniform scan and
then reports the shortest closed order-statistic window at that count. Across
the feasible cells of the prespecified simulation panel, the frozen procedure
shows reliable empirical content attainment and substantially less width
inflation than full-range Wilks intervals in many settings. That evidence is
repeated-sampling validation, not the missing exact recursion or a finite-
sample theorem for the adaptive closed-window action. The pharmaceutical study
illustrates placement and held-out stability on historical batches; it is not
a regulatory validation study. The regression and dynamic endpoint sections
extend the target algebra, but without matching software and experiments they
remain proposed methods.

\section{Cross-cutting lessons}
\label{sec:synthesis-lessons}

\subsection{The target determines the meaning of uncertainty}

The most general conclusion is that uncertainty cannot be interpreted from
computational form alone. In the dynamic quantile chapters, posterior draws may
describe model parameters, latent quantile states, or future responses after a
defined predictive step. In the MTI chapter, generalized-posterior draws
describe endpoints for a fixed loss-defined target. A tolerance statement adds
a repeated-sampling guarantee for the reported interval rule. Moving between
these levels requires an explicit probabilistic construction, calibration
argument, or validation study.

This distinction also governs interval terminology. A credible band for a
fitted quantile curve is not a predictive band for a future response. A pair of
quantile estimates does not represent every fixed-content interval target. A
high posterior probability that selected endpoints have a property does not by
itself establish tolerance confidence. Keeping these statements separate is
not merely terminological; it prevents evidence for one object from being used
to support a different claim.

\subsection{Approximation should be evaluated as part of the method}

All four project lines use computation to expose otherwise difficult targets.
MCMC provides a reference posterior calculation in several models, while VB or
LDVB reduces cost by optimizing a restricted family. Gaussian augmentation
enables state simulation and coefficient updates. Numerical scan calibration
defines the TCSP action. These devices are substantive parts of the procedures,
not neutral implementation details.

Consequently, approximation diagnostics must follow the claim. ELBO behavior,
iteration limits, and comparisons with MCMC assess aspects of variational
computation. Score and calibration summaries assess fitted or predictive
behavior. Uniform-scan calibration and independent simulation assess the
tolerance action. No single diagnostic validates all three levels. The mixed
results preserved in the research chapters are useful precisely because they
show where speed, uncertainty, and predictive performance do not move together.

\subsection{Provenance is part of empirical meaning}

The application results depend on source versions, forecast origins, horizon
definitions, covariate availability, reservoir selections, and evaluation
protocols. The dissertation therefore records exact manuscript and package
snapshots and treats the external research repositories as the computational
authority. Selected tables and figures are readable within the thesis, but
code, data, fitted objects, and full output archives are not duplicated here.

This separation avoids turning the dissertation into a second, divergent
research monorepository. It also narrows reproducibility claims. The thesis
build is self-contained and the imported scientific artifacts have recorded
source hashes. That establishes document provenance and technical buildability;
it does not establish bitwise regeneration of every scientific result or grant
permission to redistribute every source asset outside the approved setting.

\section{Limitations shared across the dissertation}
\label{sec:synthesis-limitations}

Several limitations recur across project boundaries.

\begin{enumerate}[label=(\arabic*)]
\item \textit{Working-likelihood calibration.} AL and exAL formulations provide
tractable conditional-quantile analyses, but posterior spread under
misspecification requires direct assessment. Sandwich adjustments, learning-
rate calibration, and other generalized-Bayes corrections remain relevant.

\item \textit{Dependence across quantile levels.} Separate fits plus synthesis
or rearrangement do not define a joint posterior over quantile-specific
parameters. The Q--DESN joint formulation shares information but does not
guarantee noncrossing. Richer joint models must balance coherence, scale, and
computational cost.

\item \textit{Forecast-design breadth.} The hydrologic study uses five origins,
and the Q--DESN GloFAS study uses one retrospective origin. The PriceFM design
uses three folds but depends on region-specific selection and retrospectively
observed leads. Broader claims require multi-origin studies with input
information available at each historical forecast origin.

\item \textit{Approximation and feature uncertainty.} The VB methods require
more systematic calibration against posterior simulation, especially for
tails and interval widths. The Q--DESN analyses condition on selected fixed
reservoirs; variability due to reservoir construction and selection is not
integrated into the reported posterior intervals.

\item \textit{Tolerance theory and endpoint extensions.} The TCSP evidence does
not yet include an exact action-matched finite-sample proof. Regression and
dynamic MTI formulations need implementations, numerical diagnostics,
simulation designs, and target-specific tolerance calibration before empirical
performance or guarantees can be claimed.

\item \textit{External validity and source stewardship.} Each empirical result
is tied to a bounded data set and protocol. Final coauthor or publisher reuse
review remains separate from scientific validity, and granular individual
contribution allocations remain to be confirmed before final submission.
\end{enumerate}

These limitations do not collapse the four contributions into an unfinished
program. They identify the boundary between what the completed studies support
and what a subsequent study would need to establish.

\section{Directions for future research}
\label{sec:synthesis-future}

One direction is a joint dynamic framework that combines source-aware
discrepancy correction, nonlinear recurrent features, and coherent dependence
across quantile levels. Such a model should be judged against simpler separate-
quantile fits, include hard or probabilistic monotonicity controls, and report
the computational cost of the joint structure. Sparse precision methods and
targeted covariance calculations would be important for scalable MCMC and VB.

A second direction is calibration under working likelihoods. Direct comparison
of posterior intervals with repeated-sampling quantile coverage could clarify
when AL or exAL posterior spread is adequate and when sandwich or learning-rate
adjustment is needed. For variational inference, controlled paired fits using
the same model specification would separate approximation error from model-
selection differences.

A third direction is operational hydrologic evaluation. Version-consistent
forecast-origin archives would support a dense hindcast, and a hierarchical
model for future precipitation and soil-water ensembles would propagate
exogenous uncertainty rather than condition on summaries. Spatial coupling
across gauges or basins could then test whether source-specific discrepancy
learning transfers beyond a single site.

A fourth direction is a complete inferential theory for adaptive fixed-content
interval actions. For TCSP, this includes an exact scan calculation or a proof
matched to the closed-window rule. For conditional and dynamic intervals, it
requires distinguishing estimation of a loss-defined endpoint function from
content calibration and repeated-sampling tolerance confidence. Implementing
the proposed MTI regression and state-space algorithms is therefore only the
first step; each target needs a corresponding simulation and diagnostic design.

Finally, the software and provenance work suggests a methodological practice
for collaborative statistical research: versions, inferential targets, and
evidence levels should be recorded together. A reproducible command can still
answer the wrong statistical question if it regenerates an artifact whose
target or information set is unclear. Conversely, a carefully delimited result
can remain scientifically useful even when full raw-data reconstruction is not
available, provided that the limitation is visible and the supporting snapshot
is fixed.

\section{Concluding perspective}
\label{sec:synthesis-conclusion}

The projects in this dissertation share a concern with distributional features
that must adapt to time, nonlinear dependence, heterogeneous information, or
interval-placement requirements. Their combined contribution is not a single
universal model. It is a set of statistically explicit routes from a decision-
relevant target to a computational procedure and an evidence-matched claim.

For conditional quantiles, dynamic state-space models and fixed recurrent
features provide two representations of temporal structure, while software and
source-aware synthesis make those representations usable in analysis. For
fixed-content intervals, retained-mean geometry identifies placement and shows
why endpoint uncertainty and tolerance confidence must be treated separately.
Across both settings, the central principle is the same: define the statistical
object before interpreting the computation, and validate the procedure at the
level of the claim that will be made. That principle both unifies the four
research chapters and fixes the boundary of their conclusions.
